\section{Experimental evaluation}\label{sec:evaluation}

\paragraph{Status.} This section reports three studies, all regenerated
byte-for-byte by the committed scripts: a seeded synthetic sample with
planted logistic ground truth ($n=2{,}000$) that validates the harness,
and two \emph{real} datasets --- the ULB credit-card fraud
data~\cite{dalpozzolo2015} (284{,}807 transactions, 492 frauds, verified
against its published invariants) and the official replication data of
Bao et al.~\cite{bao2020} (146{,}045 firm-years, 964 AAER frauds,
fiscal 1990--2014, obtained from the authors' public repository).

\subsection{Setup}

Rows are split 70/30 with a fixed seed; quantile thresholds are fitted
on the training split only and map each feature to
\texttt{\_high}/\texttt{\_low} propositions, with a high-amount context
variable. Candidate rules (single and pair antecedents) are selected by
empirical precision with support floors; rule probabilities are the
empirical precisions --- measured, never invented. Metrics are ROC AUC,
Brier score, log-loss, and expected calibration error, implemented in
the package and unit-tested against hand-computed values.

\subsection{Scenario suite}

Table~\ref{tab:suite} reports scenario sizes and inferred confidences
for five representative bundled scenarios, generated by the benchmark
script and embedded verbatim in the working draft, where a marker-based
check enforces byte equality with the script's output.

\begin{table}[t]
\centering
\caption{Bundled scenario suite (context set ``1'' active); generated by
\texttt{scripts/benchmark.py --paper-table}.}
\label{tab:suite}
\begin{tabular}{lrrlr}
\toprule
Scenario & Facts & Rules & Query & Confidence\\
\midrule
accounting\_very\_complex & 6 & 5 & InvestigationRequired & 0.892294\\
auditing\_complex & 4 & 3 & InvestigationRequired & 0.931000\\
pharmaceutical\_very\_complex & 6 & 7 & RecallRequired & 0.995500\\
pharmaceutical\_very\_complex & 6 & 7 & PublicNotification & 0.992130\\
pharmaceutical\_very\_complex & 6 & 7 & RegulatoryInvestigation & 0.950000\\
oncology\_complex & 4 & 3 & BiopsyRequired & 0.859500\\
logistics\_very\_complex & 6 & 5 & CustomerNotification & 0.893000\\
\bottomrule
\end{tabular}
\end{table}

\subsection{Synthetic harness validation}

Table~\ref{tab:headtohead} compares \pla{} with static
empirical-precision rules (noisy-OR aggregation, neutral context),
\pla{} with learned weights (Section~\ref{sec:learning}, intercept
included), and four baselines: logistic regression and gradient boosting
on the raw numeric features, the same \pla{} rules executed as a ProbLog
program (independent rules $=$ noisy-OR; deterministic test subsample),
and a pgmpy naive Bayes over the propositions.

\begin{table}[t]
\centering
\caption{Synthetic development sample (30\% test); generated by
\texttt{scripts/run\_experiments.py} and reproduced byte-for-byte from
scratch.}
\label{tab:headtohead}
\begin{tabular}{lrrrr}
\toprule
Model & AUC & Brier & log-loss & ECE\\
\midrule
logistic regression (raw features) & 0.9746 & 0.0413 & 0.1390 & 0.0298\\
gradient boosting (raw features) & 0.9171 & 0.0883 & 0.3750 & 0.0808\\
ProbLog on \pla{} rules (subsample) & 0.7416 & 0.1327 & 1.4447 & 0.0704\\
pgmpy naive Bayes (propositions) & 0.7228 & 0.0993 & 0.3405 & 0.0288\\
\pla{} static (empirical precisions) & 0.7217 & 0.1138 & 1.0908 & 0.0952\\
\pla{} learned (logit path + bias) & 0.7227 & 0.0999 & 0.3420 & 0.0290\\
\bottomrule
\end{tabular}
\end{table}

Two observations. First, learning matches the static rules' ranking (AUC
$0.7227$ vs.\ $0.7217$) while winning calibration decisively (log-loss
$0.342$ vs.\ $1.091$): raw empirical precisions are badly overconfident
when several rules fire together, and maximum likelihood repairs
exactly that. Second, both \pla{} variants trail the raw-feature
baselines, as expected --- the synthetic ground truth is logistic in
exactly those continuous features, and quantile discretization discards
most of that signal; the informative comparisons are on real data.
The ProbLog row executes the same rules under distribution semantics and
lands near \pla{} static, cross-validating the noisy-OR implementation
across systems.

\subsection{Real data I: credit-card fraud}

Table~\ref{tab:creditcard} reports the ULB dataset (random 70/30 split,
thresholds fitted on train only; 199{,}364 train / 85{,}443 test).

\begin{table}[t]
\centering
\caption{ULB credit-card fraud (real data; fraud rate 0.17\%); generated
by \texttt{scripts/run\_experiments.py --data creditcard.csv}.}
\label{tab:creditcard}
\begin{tabular}{lrrrr}
\toprule
Model & AUC & Brier & log-loss & ECE\\
\midrule
logistic regression (raw features) & 0.9796 & 0.0006 & 0.0036 & 0.0003\\
\pla{} static (empirical precisions) & 0.9687 & 0.0018 & 0.0172 & 0.0138\\
\pla{} learned (logit path + bias) & 0.9625 & 0.0017 & 0.0113 & 0.0022\\
ProbLog on \pla{} rules ($n{=}150$) & 0.961 & 0.0439 & 0.1502 & 0.0553\\
pgmpy naive Bayes ($n{=}2000$) & 0.9322 & 0.0016 & 0.0195 & 0.0016\\
gradient boosting (untuned defaults) & 0.7159 & 0.0013 & 0.0262 & 0.0009\\
\bottomrule
\end{tabular}
\end{table}

The headline: an 11-rule, fully hand-traceable \pla{} model reaches AUC
$0.96$--$0.97$ on real fraud data, within $0.011$--$0.017$ of raw-feature
logistic regression, and the learned variant is nearly perfectly
calibrated (ECE $0.0022$). This study also surfaced a genuine
model-class pathology, preserved in the repository's history: without
the intercept, the learned variant's AUC was $0.076$ --- \emph{inverted}
--- for the reason given in Section~\ref{sec:learning}; the intercept
restores it to $0.9625$. Untuned default gradient boosting
underperforms badly at this class imbalance; we deliberately report
defaults rather than tuned baselines and note this as a limitation.

\subsection{Real data II: accounting fraud under temporal drift}

Table~\ref{tab:bao} reports the Bao et al.\ replication
data~\cite{bao2020} under a temporal split (train fiscal
$\le 2001$, gap year 2002, test $\ge 2003$; complete-case on the paper's
financial ratios: 62{,}689 train / 58{,}405 test; 13 discretized ratios,
issuance indicator as the context variable).

\begin{table}[t]
\centering
\caption{Bao et al.\ accounting fraud (real data; temporal split);
generated by \texttt{scripts/run\_experiments.py --schema bao2020}.}
\label{tab:bao}
\begin{tabular}{lrrrr}
\toprule
Model & AUC & Brier & log-loss & ECE\\
\midrule
logistic regression (14 ratios) & 0.6771 & 0.0061 & 0.0360 & 0.0006\\
gradient boosting (untuned defaults) & 0.6349 & 0.0064 & 0.0384 & 0.0011\\
\pla{} learned (logit path + bias) & 0.5343 & 0.0061 & 0.0375 & 0.0025\\
\pla{} static (empirical precisions) & 0.5312 & 0.0063 & 0.1424 & 0.0031\\
\bottomrule
\end{tabular}
\end{table}

This is an honest negative result on a deliberately simplified protocol:
quantile tail-rules mined before 2002 barely rank post-2003 frauds above
chance (AUC $\approx 0.53$), and the issuance context variable does not
rescue them --- while raw-ratio logistic regression reaches $0.677$.
Temporal drift in accounting fraud is exactly the regime the
context-conditioning hypothesis targets, and on this first, simplified
attempt the mechanism did not earn its keep; per the pre-registered kill
criterion we report that plainly. Three protocol gaps temper the
conclusion and define the next iteration: we drop the 13.4\% of
firm-years with missing ratios, we do not implement the original
paper's serial-fraud handling, and the rule vocabulary (11 tail rules)
is far coarser than the 28-feature ensembles of~\cite{bao2020}.

\subsection{Trace fidelity}

Explanation quality is measured, not asserted, with deletion metrics in
the faithfulness framing of~\cite{jacovi2020}: \emph{comprehensiveness}
(drop in prediction when the trace's top-ranked rule's facts are
removed) and \emph{sufficiency} (gap when only they are kept), against a
reversed-attribution control that ranks the same fired rules worst
first.

\begin{table}[t]
\centering
\caption{Trace fidelity (comprehensiveness, higher is better, vs.\ the
reversed control); generated by \texttt{scripts/run\_fidelity.py}.}
\label{tab:fidelity}
\begin{tabular}{llrrr}
\toprule
Dataset & Model & $n$ & Trace & Reversed control\\
\midrule
synthetic & \pla{} static & 350 & 0.2222 & 0.1864\\
synthetic & \pla{} learned & 350 & 0.0939 & 0.0807\\
credit card & \pla{} static & 48{,}630 & 0.0202 & 0.0166\\
credit card & \pla{} learned & 48{,}630 & 0.0002 & 0.0001\\
accounting & \pla{} static & 21{,}668 & 0.0163 & 0.0127\\
accounting & \pla{} learned & 21{,}668 & 0.0003 & $-0.0004$\\
\bottomrule
\end{tabular}
\end{table}

The static model's traces out-score the reversed control on every
dataset (Table~\ref{tab:fidelity}). For the calibrated learned model on
rare-event data the deletion deltas compress toward zero: predictions
live near the base rate, so removing a rule moves them by fractions of a
percentage point in absolute probability. The ordering still favors the
faithful trace, but the margins are at noise level --- a measurement
lesson we take forward: deletion metrics should be computed in log-odds
rather than probability units for calibrated rare-event models.
Hand-computed unit tests pin the metric definitions themselves.
